% The state update of Section III-D (models.v, state_step), in the notation of the two
% kinds of step, and what the two policies of Section IV put in its open stages.
% model.tex inputs this file before fig-ordinary, so that the two ordinary figures and
% the two sliced ones pair.
\begin{figure}[t]
\centering
{\footnotesize
\begin{align*}
&\stepI{v}{i}{\mathsf{set\_pending}\;i\;v} \\
&\stepO{v}{o}{\mathsf{initiate\_next}\;o\;(\mathsf{advance\_policy}\;
  (\mathsf{observe\_policy}\;o\;v))}
\end{align*}
\vspace{-1ex}
\begin{tabular}{@{}lll@{}}
\toprule
 & \code{observe_policy} & \code{advance_policy} \\
\midrule
ordinary & \code{observe_handler_output} & \code{id} \\
sliced & \code{observe_timer_output} & \code{tick_slice} \\
\bottomrule
\end{tabular}
}
\caption{How a step changes the state $v$ (\m{state_step}). An input only records the
  interrupt. An output first passes through the two stages of the service policy, the
  first of which reads the output and the second only the state, and then
  \m{initiate_next} chooses who runs next. Only the sliced policy uses the second
  stage.}
\label{fig:update}
\end{figure}
